% Generated draft: requires mathematical and source review.
\documentclass[11pt]{article}
\usepackage[margin=25mm]{geometry}
\usepackage{fontspec}
\setmainfont{Latin Modern Roman}
\usepackage{amsmath,amssymb,mathtools,unicode-math}
\setmathfont{Latin Modern Math}
\usepackage{enumitem,needspace,xurl,hyperref}
\setlength{\parindent}{0pt}
\setlength{\parskip}{5pt}
\setlength{\emergencystretch}{3em}
\newcommand{\N}{\mathbb{N}}
\newcommand{\Z}{\mathbb{Z}}
\newcommand{\Q}{\mathbb{Q}}
\newcommand{\R}{\mathbb{R}}
\newcommand{\C}{\mathbb{C}}
\newcommand{\F}{\mathbb{F}}
\newcommand{\A}{\mathbb{A}}
\newcommand{\PP}{\mathbb{P}}
\newcommand{\E}{\mathbb{E}}
\newcommand{\eps}{\varepsilon}
\newcommand{\dd}{\,\mathrm{d}}
\begin{document}
\section*{UnsolvedMath ID: 20002492}
\textbf{Source problem code:} catalogue entry
\subsection{Definitions and mathematical statement}
\textbf{Problem 36 (Massimiliano Berti).} For the autonomous massless wave equation with Dirichlet boundary conditions,
\[
 u_{tt}-u_{xx}=f(u),\qquad f(0)=f'(0)=0,
\]
find small-amplitude periodic orbits with large periods and genuinely quasiperiodic solutions. The source mentions many fixed-period small-amplitude solutions for general nonlinearities, but supplies no interval, smoothness, nonflatness/twist hypotheses, amplitude norm or precise quantifiers. The aid uses the interval \((0,L)\), \(L>0\), and \(u(t,0)=u(t,L)=0\). A meaningful ``large period'' target must concern minimal periods tending to infinity, since multiples of a fixed period are already arbitrarily large periods.

\textbf{Reported flat-germ obstruction.} Let \(f\in C^\infty(\R)\) vanish on a neighborhood \([-\delta,\delta]\) of zero, \(\delta>0\). Suppose a global finite-energy solution
\[
 (u,u_t)\in C(\R;H_0^1(0,L)\times L^2(0,L))
\]
is uniformly small enough that \(|u(t,x)|\leq\delta\) for all \(t,x\). Then it solves the linear Dirichlet wave equation. The aid reports that it is \(2L\)-periodic, and every nonzero such solution has sine expansion
\[
 u(t,x)=\sum_{j\geq1}\left[a_j\cos\left(\frac{j\pi t}{L}\right)
 +b_j\sin\left(\frac{j\pi t}{L}\right)\right]
 \sin\left(\frac{j\pi x}{L}\right)
\]
in the energy sense. If \(S=\{j\geq1:(a_j,b_j)\ne(0,0)\}\), its exact minimal period is
\[
 T_{\min}=\frac{2L}{\gcd(S)}\leq2L.
\]
The associated minimal temporal Bohr frequency module is
\[
 \frac{\pi\gcd(S)}{L}\Z,
\]
which has rank one. The zero solution is excluded from the minimal-period assertion.

Here genuine quasiperiodicity means effective temporal frequency-module rank at least two, not a periodic function represented on a larger redundant parameter torus. Consequently no uniformly small solution in this flat-germ class has arbitrarily large minimal period or genuine quasiperiodicity. The assumption is vanishing on a neighborhood, not just all derivatives vanishing at the single point zero.

Thus the literal conditions \(f(0)=f'(0)=0\) cannot guarantee the requested behavior uniformly for every smooth nonlinearity. A positive theorem must exclude this class, for example by an appropriate nonflatness or twist assumption; the aid supplies no sufficient general condition. The example does not refute results for a specified nonflat nonlinearity or a generic class.

The aid reports substantial known fixed-period regimes, while saying that exact autonomous massless Dirichlet extensions were not verified in its literature search. No primary references, dates or search scope are supplied for that claim. It is not a new nonexistence theorem or a determination of the present literature status.
\subsection{Short English statement}
The wave-equation question needs a nonflatness assumption and a minimal-period convention. If the nonlinearity vanishes near zero, every uniformly small finite-energy solution is linear, with minimal period at most \(2L\) and only one effective temporal frequency. This obstructs a claim uniform over all nonlinearities but leaves nonflat and generic regimes separate.
\subsection{Sources}
\begin{itemize}
\item Local UnsolvedMath registry, permanent ID \(20002492\), catalogue entry, \texttt{aim-pdes-notes.json}, \texttt{notes[103]}, \emph{New connections between dynamical systems and PDE's}, Problem 36, Massimiliano Berti: \url{https://aimath.org/WWN/dynpde/dynpde.pdf}.
\item Supplied verification identifies lost leading digit 3, page artifact 10 and line-end hyphenation; the original displayed 6 is not used as a new identity. No independent website/PDF visit is claimed.
\item Supplied research summary on flat-germ linearization, exact sine-mode minimal periods and rank-one frequency module. Attempt 1: partial\_result, partially\_solved, rigor ``proved'', low novelty confidence. The machine-generated aid is not peer reviewed and must be independently verified before citation. It mentions fixed-period primary results and a literature search without supplying full references or dates.
\end{itemize}
\end{document}
